\documentclass[12pt]{article}
\usepackage[T1]{fontenc}
\usepackage[utf8]{inputenc}
\usepackage[brazil]{babel}
\usepackage[margin=1in]{geometry}
\setlength{\parindent}{0pt}
\begin{document}
\section*{Tema 5}
Processo: PEDILEF 2006.71.95.01 8143-8/ RS\\
Situacao: Revisado pelo Tema 145\\
Ramo: DIREITO PREVIDENCIÁRIO\\
Questao: Saber se para obtenção da aposentadoria por idade de segurado especial é necessário demonstrar atividade rural no período imediatamente anterior ao atingimento da idade mínima ou à apresentação do requerimento administrativo.\\
Tese: Para a obtenção de aposentadoria por idade rural, é indispensável o exercício e a demonstração da atividade campesina correspondente à carência no período imediatamente anterior ao atingimento da idade mínima ou ao requerimento administrativo (TESE FIRMADA NO TEMA 145).\\
Fonte: https://www2.jf.jus.br/phpdoc/virtus/mostradocumento.php?Process=200671950181438\&seq=3
\end{document}
